\section*{2. A finite window at the river}
Write \(V(t,y)\) for the generating function of diagrams in \(U\) with
no arch common to both matchings. A common arch separates its endpoints
from the other vertices in the noncrossing join: its interior can have
no arch going to the exterior on either side. Conversely, a two-point
join block is a common arch. Removing the singleton term \(z\) in the
block decomposition therefore gives
\[
             V=F(tV^2,y)-tV^2.                                      \tag{3}
\]
We construct an upper bound on the weight of all continuations of a
partially drawn diagram. This supplies a lower bound on the convergence
radius of \(V\).

Scan the river vertices from left to right. A vertical cut meets two
lists of arch ends, one above and one below the river. Write each list
in order of increasing distance from the river. In the scanned half-plane
the ends are joined in pairs by disjoint paths. A length counts the
visited vertices on one path. The side alternates on crossing a vertex,
so opposite-side ends have odd length and same-side ends have even
length. Record a length by class 1 for length one, 2 for positive even
length, and 3 for odd length at least three. The four possible moves,
specified by the opening or closing parenthesis in each matching, are:
\begin{itemize}
\item \(U\): open on both sides, creating a path of length one with a
new nearest end on each list.
\item \(R,L\): close on one side and open on the other. Move the nearest
end of the closing list to the head of the other list and increment
its path length.
\item \(D\): close on both sides, consuming the nearest end on each
list. Two different paths of lengths \(i,j\) merge to length \(i+j+1\).
A shared path closes with weight \(q\); closure of a class-1 path is
forbidden because it creates a common arch.
\end{itemize}
A move requiring a missing end is unavailable. The two parentheses
at each vertex prescribe the exposed ends uniquely, by noncrossing
on each side. Consequently histories from empty to empty are in
bijection with the labelled arch pairs entering \(V\).

Let \(\mathcal S\) consist of all finite noncrossing pairings of these
ends with these parity-compatible classes. This includes every physically
reachable frontier. The four rules act on \(\mathcal S\) without actual
lengths: increment takes classes \(1,2,3\) to \(2,3,2\); a merge always
has length at least three and its parity is known. For positive weights
\(w_1,w_2,w_3\) put \(H(s)=\prod_p w_{c(p)}\) over open paths of \(s\).
For a valid move \(e:s\longrightarrow es\) let \(r_e(s)\) be
\(H(es)/H(s)\) times the closing weight, if any.

\begin{lemma}
Retain the nearest \(d\) ends of each list, or the entire list if
shorter. Also retain their partners and the classes of their paths.
Erase all other paths and compress the two lists, preserving order.
Call the resulting pairing \(P_ds\). For a move \(e\), its availability,
its multiplier, and \(P_d(es)\) are determined by \(P_{d+1}s\).
\end{lemma}
Indeed, at most one existing end is removed from either list in a single
move. Any old end that becomes one of its nearest \(d\) was among the
previous nearest \(d+1\) on its side (including an end about to slide
to the other side). A new end is given explicitly. The changed paths
meet the old nearest ends. All their old partners are therefore present,
including the far ends of two paths that merge. The order among
unaffected retained ends is preserved. If the full list had fewer than
\(d+1\) ends it was retained in its entirety. These facts also prove
\(P_dP_{d+1}=P_d\). This argument holds at arbitrary frontier size.

To enumerate the possible \(P_ds\), take all noncrossing pairings
of \(2j\) ends with \(0\le j\le 2d\), split the ordered ends at
every possible divider, and discard the pairing if some path misses
both near-end windows. The linear boundary order reverses the first
list and then follows the second. Assign class 2 to each same-side
path and either 1 or 3 to each opposite-side path. Every possible
window occurs in this list. Conversely each enumerated configuration
is its own window in \(\mathcal S\). For functions symmetric in the
two lists we identify a configuration with its reflection. Both slides
are still included in the sum, even if they give the same target.

Define the positive linear operator \(Tg(s)=\sum_e r_e(s)g(es)\).
A finite test for a multiplier bound \(\Lambda\) consists of positive
numbers \(h(P_ds)\) such that for each enumerated width-\((d+1)\)
window \(S\),
\[
       \sum_e r_e(S)\,h(P_d(eS))\ \le\ \Lambda h(P_dS).                \tag{4}
\]
The same \(S\) determines all terms on the left. The lemma gives (4)
on the whole state space. Iteration from the empty state bounds the
total closing weight of histories returning empty after \(2n\) steps by
\(\Lambda^{2n}\), since their initial and final potentials coincide.
Accordingly \(V(t,1/q)\) has radius at least \(q/\Lambda^2\).

\paragraph{What increasing the width can prove.}
For fixed \(q,w_i>0\) put \(a_j=\sup_{\mathcal S} T^j1\) and
\(\rho=\lim_{j\to\infty}a_j^{1/j}\). This limit exists by
submultiplicativity; each \(a_j\) is finite since the local multipliers
are bounded. The infima of the constants in (4) decrease with \(d\)
and converge to \(\rho\). In one direction any positive finite-window
potential has positive minimum and finite maximum. Iteration of (4)
thus gives \(a_j\le (\max h/\min h)\Lambda^j\). For the other direction
take \(\lambda>\rho\) and choose \(L\) such that \(a_L<\lambda^L\).
Each \(T^j1\) depends only on \(P_j s\), by the lemma. The positive
reflection-invariant function
\[
 g=\sum_{j=0}^{L-1}\lambda^{-j}T^j1
 \qquad\hbox{satisfies}\qquad
 Tg=\lambda(g-1+\lambda^{-L}T^L1)\le\lambda g .
\]
It is a width-\((L-1)\) certificate. The monotonicity uses
\(P_dP_{d+1}=P_d\). Here \(\rho\) is the uniform growth of continuations
from arbitrary formal states. Identifying it with the empty-to-empty
growth rate is a separate question.

\paragraph{The rational certificate.}
Take \(q_0=18\), \(d=6\), \((w_1,w_2,w_3)=(1432,2940,4284)/1000\),
and \(\Lambda=5447/1000\). Up to reflection there are 25,000 width-six
windows and 128,600 width-seven extensions. The certificate enumerates
them by matching the first end of each interval to an admissible partner
and splitting into its interior and exterior subintervals. A pair with
both endpoints outside the retained window is excluded at its creation.
A separate count by reflection and Burnside's lemma verifies both totals.
Start every width-six entry at \(10^{10}\). For each source take the
maximum over its extensions of the left side of (4). Normalize the new
vector by its maximum, multiply by \(10^{10}\), round down and replace
entries below one by one. After 180 iterations the largest exact ratio
over the extensions is
\[
   299370707754862/55041444031875 < 5.447 .
\]
Only the final inequalities are needed for the bound.

Set \(t_0=6065/10000\) and \(b=6977/10000\); then
\(t_0\Lambda^2<18\). For a lower bound let \(K_a\) be the sum of
\(18^{\#\text{loops}}\) over empty-to-empty histories of length \(2a\)
with at most 14 active ends at every stage and no common arch.
The forward recurrence starts at the empty state with weight one.
It drops moves exceeding the cap, assigns multiplier 18 to an allowed
closure, and sums incoming weights at the next state. All other
multipliers are one. It visits 18,924 reflected states and gives
\[
        t_0\left(\sum_{a=0}^{40}K_a(t_0/18)^a\right)^2 > b .
\]
Thus \(z_0=t_0V(t_0,1/18)^2>b\), and (3) implies convergence of
\(F(z,1/18)\) beyond \(b\). As a second check, differentiating (3)
at \(z=tV^2\) gives
\[
 (2zF_z-F)(V+2tV_t)=V^2(tV-1+2t^2V_t).
\]
The same restricted counts give a lower bound greater than 0.00157 for
\(t_0V-1+2t_0^2V_t\). The derivatives exist at this point because
\(t_0\) is strictly inside the radius and \(tV(t,1/18)^2\) is increasing.

For \(q=18x^2\), \(x\ge1\), the support inequality \(2s\ge a\) for all
nontrivial blocks now implies convergence of \(F(z,1/q)\) beyond \(bx\).
At that argument a lower bound for \(2zF_z-F\) is
\[
 bx-1+\sum_{2\le a\le12,s} C_{a,s}x^{a-2s},\qquad
 C_{a,s}=(2a-1)f_{a,s} b^a18^{-s}.
\]
The coefficients through order twelve follow by reverting (1) from the
complete loop-count rows. Exact arithmetic gives a value greater than
0.00016 at \(x=1\) and derivative at least
\(b-\sum(2s-a)C_{a,s}>0.4948\) throughout \(x\ge1\).
The mean \(zF_z/F\) starts at zero and has strictly positive derivative
\(\Var A\) with respect to \(\log z\). This establishes the interior
saddle throughout the claimed interval. At positive \(y\), convergence
extends to a neighbourhood in both variables since \(s\le a\).
Near \((z,y)=(1,0)\), the bounds \(f_{a,s}\le16^a\) and \(s\ge a/2\)
give an analytic neighbourhood. The saddle there is analytic by the
implicit function theorem applied to \(2zF_z-F\).
